\section{Proof of the Measure Equivalence Theorem}
\label{sec:proof-of-the-measure-equivalence-theorem}

We now assemble the finite-dimensional coercivity and the measurable
zero--one law.
Only the local coercivity theorem from \cref{sec:finite-dimensional-hellinger-coercivity} is used
in the necessity argument.
The global finite-dimensional estimate is proved afterward.

\subsection{Sufficiency of Hilbert--Schmidt perturbations}
\label{sec:sufficiency}

\begin{proposition}[Hilbert--Schmidt sufficiency and likelihood ratio convergence]
\label{prop:hilbert-schmidt-sufficiency}
Assume \eqref{eq:fisher-regularity}, let
$\mu=\rho^{\otimes\N}$, where $\rho(dx)=f(x)\,dx$, and let
$T=I+K\in\GL(H)$ with $K\in\mathcal S_2$.
Then \(\nu_T\sim\mu\).
For the eventually invertible operators
$T_n=I+P_nKP_n$, their likelihood ratios
$Z_{T_n}=d\nu_{T_n}/d\mu$ converge in $L^1(\mu)$ to the strictly
positive likelihood ratio $Z_T=d\nu_T/d\mu$, and
$\sqrt{Z_{T_n}}\to\sqrt{Z_T}$ in $L^2(\mu)$.
The inverse approximants satisfy the same conclusion:
\[
 \frac{d\nu_{T_n^{-1}}}{d\mu}
 \longrightarrow\frac{d\nu_{T^{-1}}}{d\mu}
 \quad\text{in }L^1(\mu).
\]
For $S,T\in\GL(H)$ with $S-I,T-I\in\mathcal S_2$, the
likelihood ratios obey, on the common measurable carrier of their countable
operator group,
\begin{align}
 Z_{ST}(y)&=Z_S(y)Z_T(\Phi_{S^{-1}}y),
   \label{eq:rn-cocycle}\\
 Z_T(y)Z_{T^{-1}}(\Phi_{T^{-1}}y)&=1.
   \label{eq:rn-inverse}
\end{align}
If $Q$ is an exact linear stabilizer of $\mu$, the same equivalence
holds whenever \(TQ^{-1}-I\in\mathcal S_2\) and \(T\) is invertible.
\end{proposition}

\begin{proof}
Since $P_nKP_n\to K$ in Hilbert--Schmidt norm, $T_n\to T$ in
operator norm.
The approximants are therefore eventually invertible.
Their inverse norms are uniformly bounded.
If $A_n$ is the first $n$-coordinate block of $T_n$, then
\begin{equation}
 Z_{T_n}(x)=|\det A_n|^{-1}
     \prod_{i=1}^n\frac{f((A_n^{-1}\pi_n x)_i)}{f(x_i)}.
 \label{eq:finite-approximant-density}
\end{equation}
These densities are positive almost everywhere.
For all sufficiently large $m,n$, the operator $T_n^{-1}T_m-I$
acts in a finite coordinate block and has operator norm at most $1/2$.
Hellinger invariance under the finite invertible map $T_n^{-1}$ and
\eqref{eq:finite-hellinger-upper} give
\begin{equation}
 \|\sqrt{Z_{T_m}}-\sqrt{Z_{T_n}}\|_2
 \le\sqrt{C_f}\,\|T_n^{-1}\|_{\mathrm{op}}
                    \|T_m-T_n\|_{\mathcal S_2}.
 \label{eq:approximant-hellinger-cauchy}
\end{equation}
Thus the square roots are Cauchy in $L^2(\mu)$.
They have a nonnegative limit $r$ of norm one.
Cauchy--Schwarz also gives
\begin{equation}
 \|Z_{T_m}-Z_{T_n}\|_1\le2\|\sqrt{Z_{T_m}}-\sqrt{Z_{T_n}}\|_2,
 \qquad Z_{T_n}\longrightarrow r^2\quad\text{in }L^1(\mu).
 \label{eq:approximant-l1-cauchy}
\end{equation}
For each fixed finite set of output coordinates, the row-series
realizations $\Phi_{T_n}\mathbf X$ converge to $\Phi_T\mathbf X$ in $L^2$,
where $\mathbf X=(X_j)_{j\ge1}$ has law $\mu$.
Bounded continuous cylinder tests identify $r^2\mu$ with $\nu_T$.

The inverse approximants converge in Hilbert--Schmidt norm because
\[
 \|T_n^{-1}-T^{-1}\|_{\mathcal S_2}
 \le\|T_n^{-1}\|_{\mathrm{op}}
       \|T_n-T\|_{\mathcal S_2}\|T^{-1}\|_{\mathrm{op}}.
\]
Apply the same argument to these finite-coordinate approximants.
This gives $\nu_{T^{-1}}\ll\mu$.
The measurable row-series construction in
\cref{lem:measurable-action} requires only the centered finite
second moment.
On its countable-group carrier it gives
\[
 \mu=(\Phi_T)_\#\nu_{T^{-1}}
       \ll(\Phi_T)_\#\mu=\nu_T.
\]
This proves equivalence and strict positivity of $r^2$.
Change of variables on the same carrier proves the cocycle and inverse
identities.
Finally, $\nu_Q=\mu$ and $T=(TQ^{-1})Q$ give the stabilizer
assertion.
\end{proof}

\subsubsection{The orthogonal case requires only location Fisher information}
\label{sec:orthogonal-transformations}

For orthogonal transformations the assumptions can be reduced to
\begin{equation}
 f>0\ \text{almost everywhere},\qquad
 g=\sqrt f\in W^{1,2}(\R),\qquad
 \E X=0,\qquad\E X^2=1.
 \label{eq:orthogonal-fisher-regularity}
\end{equation}
In particular, $xg'\in L^2$ is not required.
The cutoff identities $\E s(X)=0$ and $\E[Xs(X)]=1$ still follow
from $J<\infty$ and the second moment.
For a skew-symmetric matrix $A$, the score has only off-diagonal terms and
\begin{equation}
 \E S_A^2=(J-1)\|A\|_F^2.
 \label{eq:orthogonal-fisher}
\end{equation}
The generator formula remains valid.
Only $x_j\partial_i\Omega_n$ with $i\ne j$ occur.
One can justify it by radial cutoff and radial mollification.
Those operators commute with the orthogonal one-parameter group.

\begin{proposition}[Orthogonal local bounds and Hilbert--Schmidt sufficiency]
\label{prop:orthogonal-sufficiency}
Under \eqref{eq:orthogonal-fisher-regularity}, every
orthogonal $U$ on $\ell^2$ with $U-I\in\mathcal S_2$ satisfies
$\nu_U\sim\mu$.  Its likelihood ratios have the same $L^1$ and
square-root $L^2$ convergence along finite-coordinate orthogonal
approximants, as do the inverse likelihood ratios.
If $f$ is non-Gaussian, there also exist dimension-free
$c,C,\delta>0$ such that
\begin{equation}
 c\|V-I\|_F\le h(V_\#\mu_n,\mu_n)
                   \le C\|V-I\|_F
 \quad(V\in O(n),\ \|V-I\|_F\le\delta).
 \label{eq:orthogonal-local-hellinger}
\end{equation}
\end{proposition}

\begin{proof}
For a skew-symmetric $A$, integration along $\exp(tA)$ gives
\begin{equation}
 h((\exp A)_\#\mu_n,\mu_n)
 \le\tfrac12\sqrt{J-1}\,\|A\|_F.
 \label{eq:orthogonal-path-upper}
\end{equation}
If $f$ is non-Gaussian, the construction of the two dual
functions $p,q$ in \cref{lem:fourier-dual-moments} uses only $s(X),X\in L^2$
and $J-1>0$.  On skew-symmetric matrices use
$H_A=\sum_{i\ne j}A_{ij}p(X_i)q(X_j)$, with no diagonal
function.  The score $S_B$ likewise has no diagonal term.
Consequently the proof of
\cref{lem:fourier-generator-bound} contains only bounded
functions times $s(X)$ or $X$, and proves the same estimate under
\eqref{eq:orthogonal-fisher-regularity}.  The half-density
pairing argument gives the local lower bound in exponential
coordinates.  For an orthogonal $V$ sufficiently close to $I$,
its power-series logarithm is real and skew-symmetric, with
$\|\log V\|_F$ comparable to $\|V-I\|_F$.  This proves
\eqref{eq:orthogonal-local-hellinger}.

For completeness, finite-coordinate orthogonal approximants
can be constructed without applying a nonorthogonal path
estimate.  Set $T_n=I+P_n(U-I)P_n$ and, for large $n$, take its
polar factor
\[
 V_n=T_n(T_n^*T_n)^{-1/2}.
\]
It is orthogonal and acts identically outside the first $n$
coordinates.  Since $T_n\to U$ in Hilbert--Schmidt norm,
\[
 \|T_n^*T_n-I\|_{\mathcal S_2}
 \le(\|T_n\|_{\mathrm{op}}+1)
                         \|T_n-U\|_{\mathcal S_2}\longrightarrow0.
\]
The smallest singular values of $T_n$ are bounded below by
some $a>0$.  Functional calculus for $|T_n|$ then gives
\[
 \|V_n-T_n\|_{\mathcal S_2}
 =\|I-|T_n|\|_{\mathcal S_2}
 \le\frac{\|T_n^*T_n-I\|_{\mathcal S_2}}{1+a}.
\]
Thus $V_n\to U$ in Hilbert--Schmidt norm.  For sufficiently
large $n,m$, $V_n^*V_m$ is close to the identity
and has a small skew-symmetric logarithm.  Applying
\eqref{eq:orthogonal-path-upper} to this logarithm proves
that the square roots of the likelihood ratios are Cauchy.  The $L^1$
limits of the likelihood ratios are identified by the same cylinder tests as in
\cref{prop:hilbert-schmidt-sufficiency}.  Apply the argument also to
$V_n^{-1}\to U^{-1}$; the measurable inverse action then
proves equivalence and the asserted convergence.
\end{proof}

\subsection{Necessity of Hilbert--Schmidt perturbations}
\label{sec:necessity}

The next lemma transports finite-dimensional local coercivity to
Hilbert--Schmidt perturbations, including finite-rank maps with infinitely
supported rows.  It is the quantitative input to the argument with swaps of tail
coordinates.

\begin{lemma}[The local bound for Hilbert--Schmidt operators]
\label{lem:local-hs-coercivity}
Let $f$ be non-Gaussian and satisfy
\eqref{eq:fisher-regularity}.  There are constants
$c,C,\delta>0$ such that every $K\in\mathcal S_2(\ell^2)$ with
$\|K\|_{\mathcal S_2}\le\delta$ satisfies
\begin{equation}
 c\|K\|_{\mathcal S_2}
 \le h(\nu_{I+K},\mu)
 \le C\|K\|_{\mathcal S_2}.
 \label{eq:local-hs-hellinger}
\end{equation}
The constants can be chosen with $\delta<1$, so $I+K$ is
automatically invertible.  Under the weaker assumptions
\eqref{eq:orthogonal-fisher-regularity}, the same conclusion
holds for orthogonal $I+K$ whenever $f$ is non-Gaussian.
This includes finite-rank perturbations whose ranges have
infinite coordinate support.
\end{lemma}

\begin{proof}
In the general linear case, take $K_n=P_nKP_n$ and apply
\eqref{eq:local-matrix-hellinger} to its finite coordinate
block.  Its untouched tail factors contribute affinity one,
so the finite-dimensional Hellinger distance is also
$h(\nu_{I+K_n},\mu)$.  The likelihood ratio convergence in
\cref{prop:hilbert-schmidt-sufficiency} gives
\[
 h(\nu_{I+K_n},\mu)
 =\left\|\sqrt{\frac{d\nu_{I+K_n}}{d\mu}}-1\right\|_2
 \longrightarrow
 \left\|\sqrt{\frac{d\nu_{I+K}}{d\mu}}-1\right\|_2
 =h(\nu_{I+K},\mu).
\]
Also $\|K_n\|_{\mathcal S_2}\to\|K\|_{\mathcal S_2}$.
Passing to the limit proves both bounds.

In the orthogonal case, use the finite-coordinate polar
approximants $V_n$ and their likelihood ratio convergence from
\cref{prop:orthogonal-sufficiency}.  They satisfy
$V_n\to I+K$ in Hilbert--Schmidt norm.  Choose the radius in
this assertion to be half the finite-dimensional radius in
\eqref{eq:orthogonal-local-hellinger}.  Then all sufficiently
large $V_n$ lie in that Frobenius ball, and
\[
 h(\nu_{V_n},\mu)\longrightarrow h(\nu_{I+K},\mu),
 \qquad
 \|V_n-I\|_{\mathcal S_2}\longrightarrow\|K\|_{\mathcal S_2}.
\]
Pass to the limit in \eqref{eq:orthogonal-local-hellinger}.
\end{proof}

\subsubsection{Characteristic-function rigidity and arbitrary mixing}
\label{sec:characteristic-function-rigidity}

Throughout this subsection, $X_j$ are independent copies of a real random
variable $X$ with law $\rho$ and
\[
  \E X=0,\qquad \E X^2=1,
\]
and $\rho$ is non-Gaussian.  The characteristic-function arguments
below require no moments of order greater than two.  Recall
\[
  \phi(t)=\E e^{itX},\qquad \mu=\rho^{\otimes\N}.
\]
The group $\mathcal P$ consists of coordinate permutations whose row
signs preserve $\rho$.  This definition applies even when $\rho$ has
no density.

\begin{lemma}[A product approximation for small coefficients]
\label{lem:characteristic-product-tail}
There is a nonnegative function $\eta$, with $\eta(\delta)\to0$ as
$\delta\downarrow0$, such that, for every finite or countable real family
$(b_j)$ with $\sum_j b_j^2<\infty$,
\begin{equation}
 \left|\prod_j\phi(tb_j)
       -\exp\!\left(-\frac{t^2}{2}\sum_jb_j^2\right)\right|
 \leq \eta(\delta)t^2\sum_jb_j^2
 \quad\text{if}\quad \sup_j|tb_j|\leq\delta.
 \label{eq:characteristic-product-tail}
\end{equation}
For a countable family, the product denotes the characteristic function
of the $L^2$ sum $\sum_jb_jX_j$.
\end{lemma}
\begin{proof}
The centered second-moment assumption gives
\[
 \phi(u)=1-\frac{u^2}{2}+o(u^2).
\]
Consequently
\[
 \eta(\delta):=
 \sup_{0<|u|\leq\delta}
 \frac{|\phi(u)-e^{-u^2/2}|}{u^2}
 \longrightarrow0,\qquad \eta(0):=0.
\]
Both $\phi(u)$ and $e^{-u^2/2}$ have modulus at most one.  Telescoping
the difference of their finite products therefore bounds it by the sum
of the differences of the factors, proving
\eqref{eq:characteristic-product-tail} for a finite family.
For a countable family, pass to the limit through finite partial sums in
$L^2$ and hence in distribution.  This argument does not take a logarithm
of the characteristic function.
\end{proof}

\begin{lemma}[Qualitative rigidity of a unit row]
\label{lem:unit-row-rigidity}
Let $a^{(n)}\in\ell^2$ satisfy $\|a^{(n)}\|_2=1$.  If
\[
 \Law\!\left(\sum_j a_j^{(n)}X_j\right)
       \Longrightarrow\rho,
\]
then
\begin{equation}
 \inf_{\substack{j\geq1,\ \varepsilon\in\{-1,1\}\\
                  \Law(\varepsilon X)=\rho}}
       \|a^{(n)}-\varepsilon e_j\|_2\longrightarrow0.
 \label{eq:unit-row-qualitative}
\end{equation}
\end{lemma}
\begin{proof}
We first describe the limits of such coefficient arrays.  Permuting
coefficients does not change the law of their sum.  Arrange each row in
nonincreasing order of absolute value, retaining its signs and padding
with zeros if necessary.  Its $j$th coefficient then has modulus at most
$j^{-1/2}$.  From any subsequence we may extract a further subsequence
on which each coefficient converges, say $a_j^{(n)}\to a_j$.
Fatou's lemma gives
\[
 q:=\sum_j a_j^2\leq1.
\]
For a fixed $m$, apply
\cref{lem:characteristic-product-tail} to the coefficients with
$j>m$.  For every fixed $t$, the error in replacing their characteristic
function by
\[
 \exp\!\left[-\frac{t^2}{2}
        \left(1-\sum_{j\leq m}(a_j^{(n)})^2\right)\right]
\]
is at most
$t^2\eta(|t|/\sqrt{m+1})$, uniformly in $n$.
Fix any $t$.
First let $n\to\infty$, and then let $m\to\infty$.
The limiting characteristic function is
\begin{equation}
 e^{-(1-q)t^2/2}\prod_j\phi(a_jt).
 \label{eq:unit-row-gaussian-summand}
\end{equation}
Equivalently, the limiting law is that of
\[
 \sum_j a_jX_j+\sqrt{1-q}\,Z,
\]
where $Z$ is standard Gaussian and independent of the $X_j$.

Under the convergence hypothesis this limiting law is the law of $X$.
Suppose that no $|a_j|$ is one.  If all $a_j$ vanish,
\eqref{eq:unit-row-gaussian-summand} already makes $X$ Gaussian.
Otherwise the maximum
$\theta:=\max_j|a_j|$ exists, since $(a_j)\in\ell^2$, and satisfies
$0<\theta<1$.  Iterating the characteristic-function identity
\eqref{eq:unit-row-gaussian-summand} $k$ times gives
\begin{equation}
 \phi(t)=e^{-(1-q^k)t^2/2}
       \prod_{\alpha\in\N^k}\phi(tb_\alpha),
 \qquad
 b_\alpha=a_{\alpha_1}\cdots a_{\alpha_k}.
 \label{eq:smoothing-iteration}
\end{equation}
Here
\[
 \sum_\alpha b_\alpha^2=q^k,
 \qquad \sup_\alpha|b_\alpha|\leq\theta^k.
\]
The identity may be obtained by substituting independent copies into
the $L^2$ series; at each finite depth all series have summable
coefficient squares.  Applying
\cref{lem:characteristic-product-tail} to the last product in
\eqref{eq:smoothing-iteration} yields
\[
 |\phi(t)-e^{-t^2/2}|
 \leq t^2q^k\eta(|t|\theta^k)\longrightarrow0.
\]
This again makes $X$ Gaussian.
This is a contradiction.
Thus some
$a_j=\varepsilon\in\{-1,1\}$; all the remaining coefficients and
the independent Gaussian summand vanish.  The limiting identity then says
$\Law(\varepsilon X)=\rho$, so the sign preserves $\rho$.

To conclude, suppose \eqref{eq:unit-row-qualitative} failed.
Take a subsequence whose distance from all the indicated unit vectors is
bounded below.
The preceding extraction produces a coefficient converging to an allowed sign
of modulus one.
Since the row norms are one, its squared distance from that signed unit
vector tends to zero:
\[
 \|a^{(n)}-\varepsilon e_j\|_2^2
       =2-2\varepsilon a_j^{(n)}\longrightarrow0.
\]
Undoing the coefficient permutations gives the required contradiction.
\end{proof}

\begin{lemma}[Uniform invariance under tail symmetries]
\label{lem:tail-symmetry-uniform}
Let $\mu$ be the i.i.d. product law and let $\nu\ll\mu$.  Let
$\mathcal T_N$ consist of the members of $\mathcal P$ that change only
finitely many coordinates and fix coordinates $1,\ldots,N$.
Then
\begin{equation}
 \sup_{S\in\mathcal T_N}h(\nu,S_\#\nu)\longrightarrow0.
 \label{eq:tail-symmetry-uniform}
\end{equation}
Also, the one-coordinate marginals of $\nu$ converge to the common
coordinate law of $\mu$, uniformly over coordinates greater than $N$.
\end{lemma}
\begin{proof}
Put $\ell=d\nu/d\mu$ and
$\ell_N=\E_\mu[\ell\mid\sigma(x_1,\ldots,x_N)]$.  The martingale
convergence theorem gives $\|\ell-\ell_N\|_{L^1(\mu)}\to0$.
Every $S\in\mathcal T_N$ preserves $\mu$ and satisfies
$\ell_N\circ S^{-1}=\ell_N$.  Thus
\[
 h(\nu,S_\#\nu)^2
 \leq\|\ell-\ell\circ S^{-1}\|_{L^1(\mu)}
 \leq2\|\ell-\ell_N\|_{L^1(\mu)},
\]
which proves \eqref{eq:tail-symmetry-uniform}.
For $i>N$, the marginal density is
$\E_\mu[\ell\mid x_i]$.  Independence and $\E_\mu\ell_N=1$ imply
$\E_\mu[\ell_N\mid x_i]=1$, so
\[
 \left\|\E_\mu[\ell\mid x_i]-1\right\|_{L^1(\mu)}
 \leq\|\ell-\ell_N\|_{L^1(\mu)}.
\]
This proves the uniform marginal assertion.
\end{proof}

\begin{lemma}[Compactness of the covariance defect]
\label{lem:compact-covariance-rigidity}
Let $T$ be bounded on $\ell^2$.  If $\nu_T\ll\mu$, then
\[
 TT^*-I\quad\text{is compact}.
\]
If $T$ is boundedly invertible, $T^*T-I$ is compact as well.
\end{lemma}
\begin{proof}
We first record uniform integrability of the squares of the linear sums
\begin{equation}
 Y_b:=\sum_jb_jX_j,\qquad \|b\|_2\leq B,
 \label{eq:uniform-linear-square-family}
\end{equation}
for each fixed $B<\infty$.  Let
\[
 X^{(M)}=X\mathbf1_{\{|X|\leq M\}}
           -\E[X\mathbf1_{\{|X|\leq M\}}],
 \qquad Y_b^{(M)}=\sum_jb_jX_j^{(M)}.
\]
The variables $X^{(M)}$ are bounded and centered, and converge to $X$
in $L^2$.  Independence gives
\[
 \sup_{\|b\|_2\leq B}\|Y_b-Y_b^{(M)}\|_2
 \leq B\|X-X^{(M)}\|_2\longrightarrow0.
\]
For fixed $M$, expanding the fourth moment of a sum of independent
centered bounded variables gives
\[
 \sup_{\|b\|_2\leq B}\E|Y_b^{(M)}|^4<\infty.
\]
The bound first holds for finite sums; it also makes their truncations
Cauchy in $L^4$, proving the countable version.  Moreover,
\[
 \sup_{\|b\|_2\leq B}
 \bigl\|Y_b^2-(Y_b^{(M)})^2\bigr\|_1
 \leq B^2\|X-X^{(M)}\|_2
                  (\|X\|_2+\|X^{(M)}\|_2)
 \longrightarrow0.
\]
Thus the squares in \eqref{eq:uniform-linear-square-family}
are uniformly approximated in $L^1$ by uniformly integrable families,
and are themselves uniformly integrable.  This uses fourth moments of
bounded truncations, not a fourth-moment assumption on $X$.

Write $\nu=\nu_T$, $\ell=d\nu/d\mu$, and
$\ell_N=\E_\mu[\ell\mid\sigma(x_1,\ldots,x_N)]$.
For a finitely supported unit vector $a$ supported on coordinates greater than
$N$, put
$F_a(x)=(\sum_i a_ix_i)^2$.  Under $\mu$, these are squares from
\eqref{eq:uniform-linear-square-family} with $B=1$; under $\nu$,
they are squares of $Y_{T^*a}$ under $\mu$, with
$\|T^*a\|_2\leq\|T\|_{\mathrm{op}}$.  They are therefore uniformly integrable
under both laws.  The truncated statistic $F_a\wedge L$ depends only
on tail coordinates.  Independence yields
\[
 \left|\E_\nu(F_a\wedge L)-\E_\mu(F_a\wedge L)\right|
 \leq L\|\ell-\ell_N\|_1.
\]
First let $N\to\infty$ with fixed $L$.
Then let $L\to\infty$.
Uniform integrability gives
\[
 \sup_{\substack{\|a\|_2=1,\ a\text{ finitely supported}\\
                  \supp a\subset\{N+1,N+2,\ldots\}}}
 \left|\langle(TT^*-I)a,a\rangle\right|\longrightarrow0.
\]
Since $TT^*-I$ is self-adjoint and bounded, density of finite vectors
in the tail subspace implies
\[
 \|(I-P_N)(TT^*-I)(I-P_N)\|_{\mathrm{op}}\longrightarrow0.
\]
The complementary rows and columns form a finite-rank operator.  Hence
$TT^*-I$ is a norm limit of finite-rank operators and is compact.
For invertible $T$, use
$T^*T-I=T^{-1}(TT^*-I)T$.
\end{proof}

\begin{lemma}[Hilbert--Schmidt necessity under a second-moment hypothesis]
\label{lem:hilbert-schmidt-necessity}
Assume that the common law has a positive density $f$, that
$g=\sqrt f\in W^{1,2}(\R)$, and that $xg'\in L^2(\R)$.
If $T\in\GL(H)$ and $\nu_T\sim\mu$, then
\begin{equation}
 T-Q\in\mathcal S_2
 \qquad\text{for some }Q\in\mathcal P.
 \label{eq:hilbert-schmidt-necessity}
\end{equation}
When $T$ is orthogonal, the condition $xg'\in L^2$ can be omitted.
\end{lemma}
\begin{proof}
By \cref{lem:compact-covariance-rigidity}, $\|T^*e_i\|_2\to1$.
By \cref{lem:tail-symmetry-uniform}, the law of
$\sum_j(T^*e_i)_jX_j$ converges to $\rho$.
Normalizing $T^*e_i$ to a unit vector changes this sum by a quantity tending to
zero in $L^2$.
Therefore \cref{lem:unit-row-rigidity} gives indices $\pi(i)$ and signs
$\varepsilon_i\in\{-1,1\}$ preserving $\rho$, outside a finite initial set, with
\begin{equation}
 u_i:=\varepsilon_i e_{\pi(i)},\qquad
 \|T^*e_i-u_i\|_2^2\longrightarrow0.
 \label{eq:qualitative-tail-matching}
 \end{equation}
The indices $\pi(i)$ are distinct sufficiently far out.
Indeed, if
$\pi(i)=\pi(j)$ for two such rows, then
$\varepsilon_iT^*e_i-\varepsilon_jT^*e_j$ has arbitrarily small norm
by \eqref{eq:qualitative-tail-matching}, whereas bounded
invertibility gives the lower bound $\sqrt2/\|T^{-1}\|_{\mathrm{op}}$.
Hence we can fix $N_0$ so that $\pi$ is injective on $i>N_0$.

There is also column approximation along this matching.
Compactness
of $T^*T-I$ and the fact that $\pi(i)$ leaves every finite set imply
\[
 \|Tu_i\|_2^2\longrightarrow1.
\]
Since $\langle Tu_i,e_i\rangle=\langle u_i,T^*e_i\rangle\to1$,
we obtain $\|Tu_i-e_i\|_2^2\to0$.  Consequently
\begin{equation}
 \eta_i:=\|T^*e_i-u_i\|_2^2+\|Tu_i-e_i\|_2^2\longrightarrow0.
 \label{eq:row-column-tail-errors}
\end{equation}

We prove that $\sum_i\eta_i<\infty$ on this tail.
By
\cref{lem:local-hs-coercivity}, there are $c_0,\delta_0>0$ such
that
\begin{equation}
 h(\mu,\nu_W)\geq c_0\|W-I\|_{\mathcal S_2}
 \quad\text{when }W\in\GL(H),\quad
 \|W-I\|_{\mathcal S_2}\leq\delta_0.
 \label{eq:local-lower-used}
\end{equation}
Suppose instead that $\sum_i\eta_i=\infty$.
Choose $\delta>0$ small enough that
$3\|T^{-1}\|_{\mathrm{op}}\delta<\delta_0$.  In any sufficiently remote tail,
where $\eta_i\leq\delta^2$, choose a finite set $E$ satisfying
\begin{equation}
 \delta^2\leq\sum_{i\in E}\eta_i\leq2\delta^2.
 \label{eq:active-tail-error}
\end{equation}
Pair each $i\in E$ with a distinct auxiliary index
$k_i$ beyond $E$, and set $G=\{k_i:i\in E\}$.
Since $\eta_i\to0$, these auxiliary indices may be chosen with
an arbitrarily small value of $\sum_{k\in G}\eta_k$, while $E$ remains fixed.

Let $S$ exchange each output pair $(e_i,e_{k_i})$, and let $P$ exchange
the aligned signed input pairs $(u_i,u_{k_i})$, acting identically on
their orthogonal complements.  Both are involutions in
$\mathcal P$, so they preserve $\mu$.  Define
\[
 W=P T^{-1} S T.
\]
Then
\begin{equation}
 W-I=P T^{-1}(ST-TP),\qquad
 \frac{\|ST-TP\|_{\mathcal S_2}}{\|T\|_{\mathrm{op}}}
 \leq\|W-I\|_{\mathcal S_2}
 \leq\|T^{-1}\|_{\mathrm{op}}\,\|ST-TP\|_{\mathcal S_2}.
 \label{eq:swap-relative-norm}
\end{equation}
These are finite-rank differences.
Indeed $S-I$ and $P-I$ have finite rank.
Hellinger invariance and $P_\#\mu=\mu$ give
\begin{equation}
 h(\mu,\nu_W)=h(\nu_T,S_\#\nu_T).
 \label{eq:swap-hellinger-identity}
\end{equation}
All maps in use lie in a countable operator group, so
\cref{lem:measurable-action} justifies this identity on a common
carrier.

We now estimate $ST-TP$ in \eqref{eq:swap-relative-norm}.
Obtain $T_0$ from $T$ by replacing each auxiliary row and column with
its matched coordinate vector:
\[
 T_0^*e_k=u_k,\qquad T_0u_k=e_k\qquad(k\in G).
\]
On the orthogonal complements of these input and output coordinates,
retain the corresponding block of $T$.  These conditions uniquely
determine $T_0$.
Orthogonal block decomposition then gives
\begin{equation}
 \|T-T_0\|_{\mathcal S_2}^2
 \leq\sum_{k\in G}
       \bigl(\|T^*e_k-u_k\|_2^2+\|Tu_k-e_k\|_2^2\bigr)
 =\sum_{k\in G}\eta_k.
 \label{eq:auxiliary-coordinate-approximation}
\end{equation}
For completeness, the left side counts the squared entries of the
auxiliary row and column errors once, while their intersection is
counted twice on the right.

Order the output subspaces as $E$, its complement outside $E\cup G$, and $G$.
Order the input subspaces as the aligned signed coordinates for $E$, their
complement outside the aligned $E\cup G$, and the aligned coordinates for
$G$.
In these orthonormal decompositions,
\[
 T_0=\begin{pmatrix}A&B&0\\ C&D&0\\0&0&I\end{pmatrix},
 \qquad
 ST_0-T_0P=
 \begin{pmatrix}
 0&-B&I-A\\
 C&0&-C\\
 A-I&B&0
 \end{pmatrix}.
\]
The blocks $B,C$ have finite-dimensional range or domain and are
Hilbert--Schmidt.  The sum of squared row and column errors on $E$ is
\[
 \sum_{i\in E}
 \bigl(\|T_0^*e_i-u_i\|_2^2+\|T_0u_i-e_i\|_2^2\bigr)
 =2\|A-I\|_F^2+\|B\|_{\mathcal S_2}^2+\|C\|_{\mathcal S_2}^2.
\]
The displayed block matrix gives the exact identity
\[
 \|ST_0-T_0P\|_{\mathcal S_2}^2
 =2\bigl(\|A-I\|_F^2+\|B\|_{\mathcal S_2}^2
                         +\|C\|_{\mathcal S_2}^2\bigr),
\]
and therefore
\begin{equation}
 \begin{aligned}
 2\|A-I\|_F^2+\|B\|_{\mathcal S_2}^2+\|C\|_{\mathcal S_2}^2
 &\leq\|ST_0-T_0P\|_{\mathcal S_2}^2\\
 &\leq2\bigl(2\|A-I\|_F^2+\|B\|_{\mathcal S_2}^2
                              +\|C\|_{\mathcal S_2}^2\bigr).
 \end{aligned}
 \label{eq:swap-block-energy}
\end{equation}
As the auxiliary error in
\eqref{eq:auxiliary-coordinate-approximation} tends to zero,
the displayed sum of squared errors tends to $\sum_{i\in E}\eta_i$.
Also,
\[
 \|(ST-TP)-(ST_0-T_0P)\|_{\mathcal S_2}
 \leq2\|T-T_0\|_{\mathcal S_2}\longrightarrow0.
\]
By \eqref{eq:active-tail-error} and
\eqref{eq:swap-block-energy}, we can choose actual finite auxiliary indices
so that
\[
 \frac{\delta}{2}\leq\|ST-TP\|_{\mathcal S_2}\leq3\delta.
\]
It follows from \eqref{eq:swap-relative-norm} that
\[
 \frac{\delta}{2\|T\|_{\mathrm{op}}}
 \leq\|W-I\|_{\mathcal S_2}
 \leq3\|T^{-1}\|_{\mathrm{op}}\delta<\delta_0.
\]
The local bound \eqref{eq:local-lower-used} now forces
$h(\mu,\nu_W)\geq c_0\delta/(2\|T\|_{\mathrm{op}})$.
On the other hand, $S$ is supported in the original remote tail.
\cref{lem:tail-symmetry-uniform} makes the right side of
\eqref{eq:swap-hellinger-identity} uniformly smaller than this positive
number when that tail is far enough out.
This is the contradiction.
Therefore $\sum_i\eta_i<\infty$.

Define $Q_0$ by the matched rows $u_i$ after the fixed initial cutoff
$N_0$, and by zero rows for $i\leq N_0$.
The summability just proved implies
\[
 \begin{aligned}
 \|T-Q_0\|_{\mathcal S_2}^2
 &=\sum_{i\leq N_0}\|T^*e_i\|_2^2
                      +\sum_{i>N_0}\|T^*e_i-u_i\|_2^2\\
 &\leq\sum_{i\leq N_0}\|T^*e_i\|_2^2+\sum_{i>N_0}\eta_i<\infty.
 \end{aligned}
 \]
Therefore $Q_0$ is a compact perturbation of an invertible operator and is
Fredholm of index zero.
Its cokernel consists precisely of the first $N_0$ output coordinates.
Its kernel is spanned by the unused input coordinates.
There are exactly $N_0$ unused columns.
Match them to the first $N_0$ rows with sign $+1$.
This completes $Q_0$ to a single $Q\in\mathcal P$ and proves
\eqref{eq:hilbert-schmidt-necessity}.

If $T$ is orthogonal, both covariance defects vanish, and every
operator $W=P T^*ST$ constructed above is orthogonal.
Consequently only the orthogonal part of
\cref{lem:local-hs-coercivity} is needed.  That part uses location
Fisher information alone, so the condition $xg'\in L^2$ may be
omitted in this case.
\end{proof}

\subsection{Exact identification of the linear stabilizer}
\label{sec:exact-stabilizer}

The equivalence theorem needs the Sobolev assumptions.
It constructs and controls likelihood ratios.
Exact invariance is more rigid.
It requires only the second moment.

\begin{proposition}[Exact linear stabilizer under a second-moment hypothesis]
\label{prop:exact-linear-stabilizer}
For a centered, variance-one, non-Gaussian i.i.d. product law,
the bounded invertible linear stabilizer is exactly $\mathcal P$.
This assertion requires neither a density nor Fisher information.
\end{proposition}
\begin{proof}
Every member of $\mathcal P$ preserves the product law.
Conversely, if $\nu_T=\mu$, covariance gives $TT^*=I$.
Since $T$ is boundedly invertible, it is orthogonal.  Each of its unit
rows gives a linear combination with the same law as $X$.
Applying \cref{lem:unit-row-rigidity} to the constant sequence
consisting of this row shows that the row is an allowed signed unit
vector.  Orthogonality makes the selected coordinates a bijection.
Thus $T\in\mathcal P$.
\end{proof}
\subsection{Two-sided uniform integrability and completion of the proof}
\label{sec:likelihood-convergence}

There are two likelihood ratio limits in the argument.
The finite-coordinate approximants in \cref{prop:hilbert-schmidt-sufficiency} converge in
\(L^1(\mu)\).
Their square roots converge in \(L^2(\mu)\).
Independently, once equivalence is known, the finite-marginal likelihood ratios and
their reciprocals are the two conditional-expectation martingales in
\eqref{eq:forward-likelihood-ui}--\eqref{eq:reverse-likelihood-ui}.
This gives uniform integrability in both directions.
It also rules out loss of mass at the projective limit.

\begin{proof}[Proof of \cref{thm:measure-equivalence}]
Necessity follows from \cref{lem:hilbert-schmidt-necessity}.  If
$T-Q\in\mathcal S_2$ for $Q\in\mathcal P$, then
$TQ^{-1}=I+(T-Q)Q^{-1}$ is invertible and differs from the identity
by a Hilbert--Schmidt operator.  By \cref{prop:hilbert-schmidt-sufficiency},
$\nu_{TQ^{-1}}\sim\mu$.  Since $Q_\#\mu=\mu$, this is
$\nu_T\sim\mu$.  The same proposition gives the finite-coordinate
likelihood ratio and square-root limits stated in the theorem.
The exact stabilizer is identified in \cref{prop:exact-linear-stabilizer}.
Finally, \cref{lem:orbit-dichotomy,lem:finite-marginal-likelihood-ratios}
give the finite-marginal affinity characterization, the singular alternative,
and both uniform-integrability statements.
\end{proof}
\begin{theorem}[Complete orthogonal criterion with only location Fisher information]
\label{thm:orthogonal-equivalence}
Let $f>0$ almost everywhere be a centered, variance-one, non-Gaussian
density, and assume only $\sqrt f\in W^{1,2}(\R)$ in addition
to its second moment.  For every orthogonal $U$ on $\ell^2$,
\[
 \nu_U\sim\mu
 \quad\Longleftrightarrow\quad
 U-Q\in\mathcal S_2
 \quad\text{for some }Q\in\mathcal P.
\]
If this condition fails, $\nu_U\perp\mu$.
\end{theorem}
\begin{proof}
Necessity is \cref{lem:hilbert-schmidt-necessity}.  For sufficiency, write
$UQ^{-1}=I+(U-Q)Q^{-1}$.  This is orthogonal and differs from the
identity by a Hilbert--Schmidt operator.  Apply
\cref{prop:orthogonal-sufficiency}, and then use
$Q_\#\mu=\mu$.  The alternative follows from
\cref{lem:orbit-dichotomy}.
\end{proof}

\subsection{Globalization of finite-dimensional coercivity}
\label{sec:globalization-proof}

We finish the proof deferred in \cref{sec:global-coercivity}.
The key point is this.
Failure of a dimension-free separation estimate would create a single
forbidden infinite-dimensional equivalence.

\begin{proof}[Proof of \cref{thm:global-coercivity}]
Fix \(0<a\le b<\infty\).  We first prove uniform separation away from
the stabilizer: for every \(\varepsilon>0\),
\begin{equation}
 \inf\left\{
 h(\mu_n,A_\#\mu_n)^2:
 \begin{array}{l}
 n\ge1,\ A\in\GL(n,\R),\\[-2pt]
 s_{\min}(A)\ge a,\ s_{\max}(A)\le b,\\[-2pt]
 d_n(A)\ge\varepsilon
 \end{array}\right\}>0.
 \label{eq:uniform-separation}
\end{equation}
Suppose otherwise.
Then there are dimensions \(n_k\) and matrices \(A_k\) in the displayed class
with
\(h(\mu_{n_k},(A_k)_\#\mu_{n_k})\to0\).  Since
\[
 -2\log\Aff(\mu_{n_k},(A_k)_\#\mu_{n_k})
 =-2\log\left(1-\tfrac12
 h(\mu_{n_k},(A_k)_\#\mu_{n_k})^2\right),
 \]
we may pass to a subsequence.
Without changing notation, assume this sum of block energies is finite.

Decompose \(H=\bigoplus_{k\ge1}\R^{n_k}\) into consecutive
coordinate blocks.
Now define
\(A=\bigoplus_{k\ge1}A_k\).  The common singular-value bounds imply
\(A\in\GL(H)\).  Each block law $(A_k)_\#\mu_{n_k}$ is
equivalent to $\mu_{n_k}$ because $f>0$ almost everywhere and $A_k$ is
invertible.  The summability of the block energies gives
\[
 \prod_{k\ge1}\Aff\bigl(\mu_{n_k},(A_k)_\#\mu_{n_k}\bigr)>0.
\]
Kakutani's product theorem~\cite{kakutani1948equivalence} therefore
gives \(\nu_A\sim\mu\).  The already proved
\cref{thm:measure-equivalence} therefore supplies \(Q\in\mathcal P\) such that
\begin{equation}
 D:=A-Q\in\mathcal S_2(H).
 \label{eq:globalization-D}
\end{equation}

Let \(\Pi_k\) be the orthogonal projection onto the \(k\)-th output block.
The orthogonality of these blocks gives
\[
 \sum_k\|\Pi_kD\|_{\mathcal S_2}^2=\|D\|_{\mathcal S_2}^2<\infty,
 \qquad\text{hence}\qquad
 \|\Pi_kD\|_{\mathcal S_2}\longrightarrow0.
 \]
If, for some output basis vector \(e_i\) in the \(k\)-th block, the signed
basis vector \(Q^*e_i\) lies outside the corresponding input block, then
\(A^*e_i\) and \(Q^*e_i\) have disjoint supports.
The lower singular value bound then yields
\begin{equation}
 \|A^*e_i-Q^*e_i\|_2^2
 =\|A_k^*e_i\|_2^2+1\ge a^2+1.
 \label{eq:row-escape-cost}
 \end{equation}
For all sufficiently large \(k\), this contradicts
\(\|\Pi_kD\|_{\mathcal S_2}^2<a^2+1\).
Hence \(Q^*\) maps every output coordinate of each sufficiently large block
into the same input block.
Its restriction is injective between two sets of cardinality \(n_k\).
Hence it is a bijection.
Denote the resulting member of \(\mathcal P_{n_k}\) by
\(Q_k\).  It follows that
\[
 d_{n_k}(A_k)
 \le\|A_k-Q_k\|_F
 =\|\Pi_kD\|_{\mathcal S_2}\longrightarrow0,
\]
contradicting \(d_{n_k}(A_k)\ge\varepsilon\).  This proves
\eqref{eq:uniform-separation}.

Let \(c_0,C_0,\delta_0\) be the constants in
\cref{prop:local-fisher-hellinger}, decreasing \(\delta_0\) so that
\(0<\delta_0<1\).  Given \(T\), choose a minimizer
\(Q\in\mathcal P_n\) in \eqref{eq:distance-stabilizer}.  Since
\(Q_\#\mu_n=\mu_n\),
\[
 h(\mu_n,T_\#\mu_n)
 =h(\mu_n,(TQ^{-1})_\#\mu_n),
 \qquad
 \|TQ^{-1}-I\|_F=d_n(T).
\]
When \(d_n(T)\le\delta_0\), the local theorem supplies the two-sided
linear bounds for \(h\).  When \(d_n(T)>\delta_0\),
\eqref{eq:uniform-separation} with
\(\varepsilon=\delta_0\) supplies the lower bound, while \(h^2\le2\)
supplies the upper bound.  Choose $c>0$ no larger than $c_0^2$ and
the positive infimum in \eqref{eq:uniform-separation} with
$\varepsilon=\delta_0$, and take
\[
 C=\max\{C_0^2,2\delta_0^{-2}\}.
\]
These constants give \eqref{eq:global-coercivity} in all regimes.
\end{proof}
